% ============================================================
% frontmatter/abstract.tex —— 中英文摘要
% 标题/关键词格式见 setup/format.tex 第 6 节；
% 正文为宋体小四(中文) / Times New Roman 小四(英文)，固定 20 磅行距。
% ============================================================

% ---------------- 中文摘要 ----------------
\clearpage
\cnabstracttitle
\addcontentsline{toc}{chapter}{摘要}    % 如需目录中显示摘要，保留此行；否则删除

% —— 中文摘要正文（示例，替换为真实内容）——
针对广域输电网常态化动态运维巡检场景，研究多机巢起降条件下的多无人机动态协同任务规划问题。
本文构建"动态任务选择 + 选择性多无人机协同任务分配（MUAS）"的两阶段分层求解框架，
并嵌入滚动时域运行机制，兼顾求解实时性与资源利用效率。

\vspace{12pt}
\noindent\textbf{关键词：}多无人机；任务分配；航迹规划；滚动时域；差分进化

% ---------------- 英文摘要 ----------------
\clearpage
% 【P2-9】去掉英文摘要页眉。依据：师兄论文 p9/p10（摘要/Abstract）实测 y>740 区域
%   完全为空 —— 既无页眉也无页脚；规范 2.4「页眉从第一章开始」。
%   规范 3.1 的"Abstract 部分用 Times New Roman"指正文用 TNR 字体，非要求加页眉。
%   （\fancypagestyle{abstracten} 定义保留在 setup/format.tex，需要时可再启用。）
\enabstracttitle
\addcontentsline{toc}{chapter}{Abstract}  % 如需目录中显示英文摘要，保留此行；否则删除

% —— English abstract（与中文摘要对应）——
% 【P2-10】师兄 p10 只有 Abstract + 正文，正文前不再重复排一遍英文题名（原与页眉重复）。
Aiming at the routine dynamic operation and maintenance inspection of wide-area
power transmission networks, this thesis studies the dynamic cooperative mission
planning of multiple unmanned aerial vehicles (UAVs) under multi-nest launch and
recovery conditions. A two-stage hierarchical framework combining dynamic task
selection and selective multi-UAV cooperative task allocation (MUAS) is proposed,
embedded in a rolling-horizon mechanism to balance solution timeliness and resource
utilization efficiency.

\vspace{12pt}
\noindent\textbf{Key Words:} multi-UAV; task allocation; path planning; rolling horizon; differential evolution

% 英文摘要页沿用前置默认 plain 样式（页脚居中罗马页码、无页眉），无需再显式切换。
